\section*{Глава 1. Анализ предметной области}
\addcontentsline{toc}{section}{Глава 1. Анализ предметной области}
\stepcounter{section}

\subsection{Обзор существующих программных продуктов}

На рынке представлено большое количество приложений для учёта личных и семейных расходов. Большинство из них — это мобильные приложения (Android/iOS). Полноценных веб-приложений, специально ориентированных на совместный семейный бюджет, на 2026 год очень мало.

Для сравнения были выбраны наиболее популярные решения: Дзен-мани, CoinKeeper, Monefy и Money Lover. В таблице ниже приведены их ключевые характеристики с точки зрения семейного использования.

\begin{table}[h]
	\centering
	\caption{Сравнение популярных приложений для учёта финансов}
	\label{tab:competitors}
	
	{\fontsize{12}{14}\selectfont
	\setlength{\tabcolsep}{6pt}
	
	\begin{tabular}{|p{3.7cm}|p{3.7cm}|p{3.5cm}|p{3.5cm}|}
	\hline
	\textbf{Характеристика} 
		& \textbf{Дзен-мани} 
		& \textbf{Monefy} 
		& \textbf{Money Lover} \\
	\hline
	
	Платформа 
		& Мобильное 
		& Мобильное 
		& Мобильное + Web \\
	\hline
	
	Совместный бюджет 
		& Полноценный 
		& Сильно ограничен 
		& Есть (в Premium) \\
	\hline
	
	Загрузка чеков 
		& Есть (QR-сканирование) 
		& Ограничено 
		& Есть (фото) \\
	\hline
	
	Веб-версия 
		& Ограниченная 
		& Отсутствует 
		& Полноценная \\
	\hline
	
	Стоимость ключевых функций 
		& Премиум-подписка (многие отчёты и прогнозы) 
		& Премиум для расширенной аналитики и синхронизации 
		& Unlimited wallets и shared budgets — только Premium \\
	\hline
	\end{tabular}
	
	}
	\end{table}


Как видно из таблицы, у каждого рассмотренного приложения есть существенные недостатки именно при семейном использовании:

\begin{itemize}
    \item \textbf{Дзен-мани} — лучший функционал совместного бюджета, но большинство важных отчётов, прогнозов и расширенной аналитики доступны только по платной подписке.
    \item \textbf{Monefy} — очень простой и интуитивный интерфейс, но совместный доступ сильно ограничен, а полноценной веб-версии нет.
    \item \textbf{Money Lover} — имеет веб-версию, однако создание нескольких общих кошельков и расширенные возможности бюджетирования также требуют платной подписки.
\end{itemize}

Таким образом, несмотря на обилие мобильных решений, на рынке практически отсутствует удобное, бесплатное (или полностью функциональное в базовой версии) веб-приложение, которое одновременно предоставляет:

\begin{itemize}
    \item полноценную совместную работу нескольких пользователей,
    \item систему ролей и разграничения прав доступа,
    \item удобную загрузку и выгрузку файлов (чеков, отчётов),
    \item глубокую агрегацию данных и наглядную статистику,
    \item доступ с любого устройства без установки приложений.
\end{itemize}

\textbf{Вывод.} Разработка собственного веб-приложения для учёта семейных расходов позволит устранить указанные недостатки существующих решений и создать удобный, бесплатный в базовой функциональности инструмент, специально ориентированный на совместное использование всей семьёй.


\subsection{Анализ программных инструментов разработки}

Для реализации веб-приложения были рассмотрены популярные технологии и стеки разработки. Особое внимание уделялось простоте освоения, скорости разработки, наличию необходимого функционала (аутентификация, работа с базой данных, загрузка файлов, визуализация) и подходящему уровню сложности для курсового проекта.

% \textbf{Анализ backend-фреймворков}

Наиболее распространёнными Python-фреймворками для создания веб-приложений в 2026 году являются Django, Flask и FastAPI.

\textbf{Django} — мощный full-stack фреймворк, включающий в себя большое количество готовых компонентов («батарейки из коробки»): встроенную систему аутентификации, ORM, админ-панель и механизмы защиты. Главное преимущество — высокая скорость разработки и хорошая безопасность. Основной недостаток — избыточная сложность и «тяжеловесность» для небольшого учебного проекта.

\textbf{FastAPI} — современный высокопроизводительный асинхронный фреймворк для создания API. Отличается отличной поддержкой типизации и автоматической генерацией документации. Преимущества: высокая скорость работы. Недостатки: избыточен для классического CRUD-приложения с серверным рендерингом шаблонов.

\textbf{Flask} — лёгкий микрофреймворк, который предоставляет только базовый функционал. Все необходимые расширения подключаются по мере необходимости. Главные преимущества — высокая гибкость, минимализм, простота изучения и отличная читаемость кода.

% \textbf{Анализ систем управления базами данных}

Для управления базой данных рассматривались две основные СУБД: SQLite и PostgreSQL.

\textbf{SQLite} — простая файловая база данных, не требующая установки сервера. Преимущества: крайне простая настройка, нулевые затраты на развертывание, высокая скорость для небольших проектов. Недостатки: слабая поддержка нескольких пользователей одновременно, ограниченные возможности по масштабированию и ограниченный функционал.

\textbf{PostgreSQL} — полноценная клиент-серверная реляционная СУБД. Преимущества: высокая производительность при работе с несколькими пользователями, возможности сложных запросов, хорошая масштабируемость и надёжность. Недостатки: более сложная настройка по сравнению с SQLite.

% \textbf{Анализ frontend-технологий}

Для создания пользовательского интерфейса сравнивались Bootstrap 5 и полноценные JavaScript-фреймворки React и Vue.

\textbf{Bootstrap 5} — CSS-фреймворк с готовыми компонентами и встроенной адаптивностью. Преимущества: очень быстрая вёрстка, минимальный порог вхождения, удобен для адаптивности. Недостатки: менее гибкий при создании сложной интерактивной логики.

\textbf{React} и \textbf{Vue} — мощные JavaScript-фреймворки для создания веб-приложений. Преимущества: высокая интерактивность, компонентный подход, отличная производительность при большом объёме данных. Недостатки: значительно более сложные в использовании, необходимость в дополнительных инструментах (Vite, Redux/Vuex и т.д.), что избыточно для курсового проекта.


После проведённого анализа был выбран следующий стек разработки:

\begin{itemize}
    \item \textbf{Backend}: Python 3.12 + Flask 3.1
    \item \textbf{ORM}: Flask-SQLAlchemy
    \item \textbf{Аутентификация и авторизация}: Flask-Login + Flask-Bcrypt
    \item \textbf{База данных}: PostgreSQL
    \item \textbf{Frontend}: HTML5, CSS3, JavaScript, Bootstrap 5.3
    \item \textbf{Визуализация данных}: Chart.js
    \item \textbf{Дополнительные инструменты}: Git, Visual Studio Code
\end{itemize}

Данный стек обеспечивает оптимальный баланс между простотой реализации, читаемостью кода и требуемым функционалом. Flask позволяет глубоко понять принципы построения веб-приложений без избыточных абстракций, что особенно важно в рамках курсовой работы. PostgreSQL выбрана с целью приближения проекта к реальным условиям разработки.



\subsection{Формулировка цели и задач работы}

\textbf{Цель работы} — разработать веб-приложение для учёта расходов и доходов семейного бюджета, обеспечивающее удобный совместный контроль финансов, прозрачность трат и наглядную аналитику для всех членов семьи.

Для достижения поставленной цели необходимо решить следующие задачи:

\begin{enumerate}
    \item Проанализировать существующие программные продукты и технологии разработки веб-приложений для управления семейными финансами.
    \item Определить требования целевой аудитории и спроектировать функциональность приложения (диаграммы вариантов использования, User Story).
    \item Разработать модель данных (ER-диаграмму, логическую и физическую схемы базы данных).
    \item Создать макеты пользовательского интерфейса (wireframes).
    \item Реализовать систему аутентификации и авторизации пользователей с разделением прав доступа по ролям.
    \item Разработать CRUD-интерфейс для основных сущностей (транзакции, категории, бюджеты).
    \item Реализовать механизмы загрузки и выгрузки файлов (чеки, отчёты) и систему агрегации данных с визуализацией статистики.
    \item Провести тестирование приложения, выполнить деплой и оформить отчётную документацию.
\end{enumerate}

В результате работы будет создано полноценное веб-приложение, соответствующее всем требованиям курсового проекта.
